% Options for packages loaded elsewhere
\PassOptionsToPackage{unicode}{hyperref}
\PassOptionsToPackage{hyphens}{url}
\documentclass[
]{article}
\usepackage{xcolor}
\usepackage{amsmath,amssymb}
\setcounter{secnumdepth}{5}
\usepackage{iftex}
\ifPDFTeX
  \usepackage[T1]{fontenc}
  \usepackage[utf8]{inputenc}
  \usepackage{textcomp} % provide euro and other symbols
\else % if luatex or xetex
  \usepackage{unicode-math} % this also loads fontspec
  \defaultfontfeatures{Scale=MatchLowercase}
  \defaultfontfeatures[\rmfamily]{Ligatures=TeX,Scale=1}
\fi
% \usepackage{lmodern}
\ifPDFTeX\else
  % xetex/luatex font selection
\fi
% Use upquote if available, for straight quotes in verbatim environments
\IfFileExists{upquote.sty}{\usepackage{upquote}}{}
\IfFileExists{microtype.sty}{% use microtype if available
  \usepackage[]{microtype}
  \UseMicrotypeSet[protrusion]{basicmath} % disable protrusion for tt fonts
}{}
\makeatletter
\@ifundefined{KOMAClassName}{% if non-KOMA class
  \IfFileExists{parskip.sty}{%
    \usepackage{parskip}
  }{% else
    \setlength{\parindent}{0pt}
    \setlength{\parskip}{6pt plus 2pt minus 1pt}}
}{% if KOMA class
  \KOMAoptions{parskip=half}}
\makeatother
\usepackage{color}
\usepackage{fancyvrb}
\newcommand{\VerbBar}{|}
\newcommand{\VERB}{\Verb[commandchars=\\\{\}]}
\DefineVerbatimEnvironment{Highlighting}{Verbatim}{commandchars=\\\{\}}
% Add ',fontsize=\small' for more characters per line
\newenvironment{Shaded}{}{}
\newcommand{\AlertTok}[1]{\textcolor[rgb]{1.00,0.00,0.00}{\textbf{#1}}}
\newcommand{\AnnotationTok}[1]{\textcolor[rgb]{0.38,0.63,0.69}{\textbf{\textit{#1}}}}
\newcommand{\AttributeTok}[1]{\textcolor[rgb]{0.49,0.56,0.16}{#1}}
\newcommand{\BaseNTok}[1]{\textcolor[rgb]{0.25,0.63,0.44}{#1}}
\newcommand{\BuiltInTok}[1]{\textcolor[rgb]{0.00,0.50,0.00}{#1}}
\newcommand{\CharTok}[1]{\textcolor[rgb]{0.25,0.44,0.63}{#1}}
\newcommand{\CommentTok}[1]{\textcolor[rgb]{0.38,0.63,0.69}{\textit{#1}}}
\newcommand{\CommentVarTok}[1]{\textcolor[rgb]{0.38,0.63,0.69}{\textbf{\textit{#1}}}}
\newcommand{\ConstantTok}[1]{\textcolor[rgb]{0.53,0.00,0.00}{#1}}
\newcommand{\ControlFlowTok}[1]{\textcolor[rgb]{0.00,0.44,0.13}{\textbf{#1}}}
\newcommand{\DataTypeTok}[1]{\textcolor[rgb]{0.56,0.13,0.00}{#1}}
\newcommand{\DecValTok}[1]{\textcolor[rgb]{0.25,0.63,0.44}{#1}}
\newcommand{\DocumentationTok}[1]{\textcolor[rgb]{0.73,0.13,0.13}{\textit{#1}}}
\newcommand{\ErrorTok}[1]{\textcolor[rgb]{1.00,0.00,0.00}{\textbf{#1}}}
\newcommand{\ExtensionTok}[1]{#1}
\newcommand{\FloatTok}[1]{\textcolor[rgb]{0.25,0.63,0.44}{#1}}
\newcommand{\FunctionTok}[1]{\textcolor[rgb]{0.02,0.16,0.49}{#1}}
\newcommand{\ImportTok}[1]{\textcolor[rgb]{0.00,0.50,0.00}{\textbf{#1}}}
\newcommand{\InformationTok}[1]{\textcolor[rgb]{0.38,0.63,0.69}{\textbf{\textit{#1}}}}
\newcommand{\KeywordTok}[1]{\textcolor[rgb]{0.00,0.44,0.13}{\textbf{#1}}}
\newcommand{\NormalTok}[1]{#1}
\newcommand{\OperatorTok}[1]{\textcolor[rgb]{0.40,0.40,0.40}{#1}}
\newcommand{\OtherTok}[1]{\textcolor[rgb]{0.00,0.44,0.13}{#1}}
\newcommand{\PreprocessorTok}[1]{\textcolor[rgb]{0.74,0.48,0.00}{#1}}
\newcommand{\RegionMarkerTok}[1]{#1}
\newcommand{\SpecialCharTok}[1]{\textcolor[rgb]{0.25,0.44,0.63}{#1}}
\newcommand{\SpecialStringTok}[1]{\textcolor[rgb]{0.73,0.40,0.53}{#1}}
\newcommand{\StringTok}[1]{\textcolor[rgb]{0.25,0.44,0.63}{#1}}
\newcommand{\VariableTok}[1]{\textcolor[rgb]{0.10,0.09,0.49}{#1}}
\newcommand{\VerbatimStringTok}[1]{\textcolor[rgb]{0.25,0.44,0.63}{#1}}
\newcommand{\WarningTok}[1]{\textcolor[rgb]{0.38,0.63,0.69}{\textbf{\textit{#1}}}}
\usepackage{longtable,booktabs,array}
\newcounter{none} % for unnumbered tables
\usepackage{calc} % for calculating minipage widths
% Correct order of tables after \paragraph or \subparagraph
\usepackage{etoolbox}
\makeatletter
\patchcmd\longtable{\par}{\if@noskipsec\mbox{}\fi\par}{}{}
\makeatother
% Allow footnotes in longtable head/foot
\IfFileExists{footnotehyper.sty}{\usepackage{footnotehyper}}{\usepackage{footnote}}
\makesavenoteenv{longtable}
\setlength{\emergencystretch}{3em} % prevent overfull lines
\providecommand{\tightlist}{%
  \setlength{\itemsep}{0pt}\setlength{\parskip}{0pt}}
\usepackage[]{natbib}
\bibliographystyle{plainnat}
\makeatletter
\@ifpackageloaded{subcaption}{}{\usepackage{subcaption}}
\@ifpackageloaded{caption}{}{\usepackage{caption}}
\captionsetup[subfigure]{margin=0.5em}
\AtBeginDocument{%
\renewcommand*\figurename{Figure}
\renewcommand*\tablename{Table}
}
\AtBeginDocument{%
\renewcommand*\listfigurename{List of Figures}
\renewcommand*\listtablename{List of Tables}
}
\newcounter{pandoccrossref@subfigures@footnote@counter}
\newenvironment{pandoccrossrefsubfigures}{%
\setcounter{pandoccrossref@subfigures@footnote@counter}{0}
\begin{figure}\centering%
\gdef\global@pandoccrossref@subfigures@footnotes{}%
\DeclareRobustCommand{\footnote}[1]{\footnotemark%
\stepcounter{pandoccrossref@subfigures@footnote@counter}%
\ifx\global@pandoccrossref@subfigures@footnotes\empty%
\gdef\global@pandoccrossref@subfigures@footnotes{{##1}}%
\else%
\g@addto@macro\global@pandoccrossref@subfigures@footnotes{, {##1}}%
\fi}}%
{\end{figure}%
\addtocounter{footnote}{-\value{pandoccrossref@subfigures@footnote@counter}}
\@for\f:=\global@pandoccrossref@subfigures@footnotes\do{\stepcounter{footnote}\footnotetext{\f}}%
\gdef\global@pandoccrossref@subfigures@footnotes{}}
\@ifpackageloaded{float}{}{\usepackage{float}}
\floatstyle{ruled}
\@ifundefined{c@chapter}{\newfloat{codelisting}{h}{lop}}{\newfloat{codelisting}{h}{lop}[chapter]}
\floatname{codelisting}{Listing}
\newcommand*\listoflistings{\listof{codelisting}{List of Listings}}
\makeatother
\usepackage{bookmark}
\IfFileExists{xurl.sty}{\usepackage{xurl}}{} % add URL line breaks if available
\urlstyle{same}
% fallback for those not using the hyperref driver hyperxmp:
\makeatletter
\@ifundefined{xmpquote}{\newcommand{\xmpquote}[1]{#1}}{}
\makeatother
\hypersetup{
  colorlinks=true,linkcolor=red,urlcolor=red,citecolor=red,anchorcolor=red,filecolor=red,
  pdfcreator={LaTeX via pandoc}}

\author{}
\date{}

\IfFileExists{seqsplit.sty}{\usepackage{seqsplit}}{\newcommand{\seqsplit}[1]{#1}}
\protected\def\breakseq#1{\seqsplit{#1}}
\protected\def\breaktt#1{\begingroup\ttfamily\seqsplit{#1}\endgroup}
\title{Pools, Rules, and Tools: A Template-Integrated Resource Architecture}
\author{Research Template Author\\\footnotesize{Active Inference Institute}\\\footnotesize{\texttt{author@research-template.org}}\\\footnotesize{\href{https://orcid.org/0000-0002-0000-0001}{ORCID: 0000-0002-0000-0001}} \\and Template Collaborator\\\footnotesize{Active Inference Institute}\\\footnotesize{\texttt{collaborator@research-template.org}}\\\footnotesize{\href{https://orcid.org/0000-0002-0000-0002}{ORCID: 0000-0002-0000-0002}}}
\date{2026-07-05}

% manuscript/preamble.md
% LaTeX preamble configuration for the Pools, Rules, and Tools manuscript.
%
% This file is passed to the rendering pipeline before the main manuscript
% sections. It configures packages, typography, cross-reference macros, and
% custom environments.
% ============================================================
% Core packages
% ============================================================
\usepackage{hyperref}
\usepackage{booktabs}
\usepackage{listings}
\usepackage{xcolor}
\usepackage{graphicx}
\usepackage{amsmath}
\usepackage{amssymb}
\usepackage{microtype}
\usepackage{cleveref}
\usepackage{float}
\usepackage{caption}
\usepackage{subcaption}
\usepackage{geometry}
\usepackage{setspace}
\usepackage{fancyhdr}
\usepackage{titlesec}
\usepackage{enumitem}
\usepackage{tcolorbox}
\usepackage{mdframed}
% ============================================================
% Page geometry
% ============================================================
% ============================================================
% Typography
% ============================================================
% ============================================================
% Header and footer
% ============================================================
\renewcommand{\headrulewidth}{0.4pt}
% ============================================================
% Code listing style
% ============================================================
% ============================================================
% Semantic macros for this manuscript
% ============================================================
\newcommand{\module}[1]{\texttt{#1}}
\newcommand{\fondname}[1]{\texttt{#1}}
\newcommand{\ruleset}[1]{\texttt{#1}}
\newcommand{\toolname}[1]{\texttt{#1}}
\newcommand{\repopath}[1]{\texttt{#1}}
\newcommand{\token}[1]{\texttt{\{\{#1\}\}}}
\newcommand{\jsonkey}[1]{\texttt{"#1"}}
\newcommand{\exitcode}[1]{\texttt{#1}}
% ============================================================
% Coloured note boxes
% ============================================================
% ============================================================
% Cross-reference setup
% ============================================================
% cleveref configuration — use \cref{} for all cross-references

\begin{document}

\begin{titlepage}
\centering
\vspace*{2cm}
{\LARGE\bfseries Pools, Rules, and Tools: A Template-Integrated Resource Architecture\par}
\vspace{0.75em}
{\large A Meta-Project Demonstrating Fonds, Rules, and Tools Integration in Research Software\par}
\vfill
\makeatletter
{\@author\par}
\vspace{1em}
{\@date\par}
\makeatother
\vfill
\end{titlepage}
\thispagestyle{empty}
\tableofcontents
\newpage


\section{Pools, Rules, and Tools: A Template-Integrated Resource
Architecture}\label{pools-rules-and-tools-a-template-integrated-resource-architecture}

\textbf{A Meta-Project Demonstrating Fonds, Rules, and Tools Integration
in Research Software}

\begin{center}\rule{0.5\linewidth}{0.5pt}\end{center}

{\def\LTcaptype{none} % do not increment counter
\begin{longtable}[]{@{}
  >{\raggedright\arraybackslash}p{(\linewidth - 2\tabcolsep) * \real{0.5000}}
  >{\raggedright\arraybackslash}p{(\linewidth - 2\tabcolsep) * \real{0.5000}}@{}}
\toprule\noalign{}
\begin{minipage}[b]{\linewidth}\raggedright
Field
\end{minipage} & \begin{minipage}[b]{\linewidth}\raggedright
Value
\end{minipage} \\
\midrule\noalign{}
\endhead
\bottomrule\noalign{}
\endlastfoot
\textbf{Authors} & Research Template Author\texorpdfstring{\textsuperscript{1}}{1}, Template Collaborator\texorpdfstring{\textsuperscript{1}}{1} \\
\textbf{Affiliation} & \texorpdfstring{\textsuperscript{1}}{1} Active Inference Institute \\
\textbf{Correspondence} & author@research-template.org \\
\textbf{Version} & 1.0.0 \\
\textbf{Date} & 2026-07-05 \\
\textbf{License} & CC-BY-4.0 \\
\textbf{Repository} &
\href{https://github.com/docxology/template}{docxology/template} \\
\textbf{DOI} & 10.5281/zenodo.template\_pools\_rules\_tools \\
\textbf{Keywords} & research software engineering, monorepo
architecture, reproducibility, fonds, governance rules, tool discovery,
graceful degradation \\
\end{longtable}
}

\begin{center}\rule{0.5\linewidth}{0.5pt}\end{center}

\subsection{Author Contributions}\label{author-contributions}

\textbf{Research Template Author}: Conceptualisation, architecture
design, module implementation (\texttt{fonds\_reader},
\texttt{rules\_applier}, \texttt{tools\_invoker}, \texttt{integration}),
manuscript writing (all sections), validation.

\textbf{Template Collaborator}: Manuscript review, test suite design,
exemplar resource authoring (fonds, rules, tools).

\subsection{Acknowledgements}\label{acknowledgements}

The authors thank the Active Inference Institute for hosting the public
template repository and providing the infrastructure within which this
exemplar was developed. The design of the three-layer resource
architecture draws inspiration from the Unix philosophy
\citep{Raymond2003art} and the enterprise application patterns
documented by \citep{Fowler2002patterns}.

\subsection{Data Availability}\label{data-availability}

All source code, configuration files, and exemplar resources described
in this paper are available in the public template repository at
\url{https://github.com/docxology/template} under the
\breaktt{projects/templates/template\_pools\_rules\_tools/} path. The
integration pipeline is fully reproducible from source using
\breaktt{uv\ run\ python\ projects/templates/template\_pools\_rules\_tools/scripts/02\_run\_integration.py}
from the repository root. Generated manuscript variables are stored in
\breaktt{output/data/manuscript\_variables.json} and injected at render
time.

\subsection{Competing Interests}\label{competing-interests}

The authors declare no competing interests.

\newpage

\section{Abstract}\label{abstract}

Research software repositories in monorepo configurations accumulate
three categories of shared resources that individual projects must
consume without re-implementing discovery logic: \textbf{data pools}
(bibliographies, contacts, datasets), \textbf{governance rules} (style
guides, coverage thresholds, citation schemas), and \textbf{executable
tools} (code executors, validators, skill invocations). Without a
canonical integration pattern, projects either duplicate discovery logic
or silently ignore resources that fail to load --- both outcomes degrade
reproducibility and collaborative cohesion
\citep{Wilson2014best, Taschuk2017ten}.

This paper presents \breaktt{template\_pools\_rules\_tools}, a
meta-project exemplar that demonstrates how a single project can
programmatically discover, validate, and exercise all three resource
categories with zero tight coupling to any specific resource instance.
The exemplar comprises four Python modules --- \texttt{fonds\_reader},
\texttt{rules\_applier}, \texttt{tools\_invoker}, and
\texttt{integration} --- plus three thin orchestration scripts and a
fully token-injected manuscript pipeline.

The architecture (fig.~\ref{fig:architecture}) separates \emph{resource
ownership} from \emph{resource consumption}. Resources live in top-level
\texttt{fonds/}, \texttt{rules/}, and \texttt{tools/} directories and
are never modified by consumers. Each resource exposes a typed manifest
(\texttt{fonds.yaml}, \texttt{rules.yaml}, \texttt{tools.yaml}) that the
corresponding reader module uses for discovery and validation. All
readers implement graceful fallbacks: they return \texttt{None} or empty
collections when a resource is absent, log a warning via the standard
library \texttt{logging} module, and allow the integration pipeline to
continue.

In a representative pipeline run, the integration demo loaded
\{\{integration.fonds\_loaded\}\} fonds, validated
\{\{integration.rules\_sets\_ok\}\} rule sets, discovered
\{\{integration.tools\_discovered\}\} tools, and processed
\{\{integration.bib\_entries\}\} bibliography entries --- all reported
as structured JSON that populates manuscript variable tokens at render
time. Tests covering the four \texttt{src/} modules achieve \texorpdfstring{\textgreater{}=}{>=}90\% line
coverage and use real file paths rather than mocks, ensuring that
reported counts are genuine.

The \breaktt{template\_pools\_rules\_tools} exemplar provides a reference
implementation that any project in the template repository can consult
when designing its own resource-consumption layer.

\textbf{Keywords:} research software engineering, monorepo architecture,
reproducibility, fonds, governance rules, tool discovery, graceful
degradation

\newpage

\section{Introduction}\label{introduction}

\subsection{Motivation}\label{motivation}

Modern research software repositories increasingly adopt monorepo
designs in which multiple projects share a common set of curated
resources. A monorepo consolidates source code, documentation, datasets,
and governance artefacts under one version-controlled root, enabling
atomic cross-project changes and a single source of truth for shared
data \citep{Fowler2002patterns}. The practical benefit is significant: a
bibliography updated once in
\breaktt{fonds/templates/template\_bibliography/} is immediately
available to every project that discovers it at runtime, without any
per-project copy.

Three categories of shared resource appear consistently across research
template repositories:

\begin{enumerate}
\def\labelenumi{\arabic{enumi}.}
\tightlist
\item
  \textbf{Data pools (fonds)}: curated reference sets ---
  bibliographies, contact registries, dataset catalogues --- that
  projects query but must never mutate.
\item
  \textbf{Governance rules}: machine-readable constraint schemas and
  human-readable style guidelines that projects load to validate their
  own outputs.
\item
  \textbf{Executable tools}: script-based entry points that projects
  invoke to run computations, validate artefacts, or call external
  agents.
\end{enumerate}

Without a canonical integration pattern for consuming these resources,
projects face a dilemma: they can hard-code discovery paths (creating
fragile, repo-root-sensitive logic) or skip resource consumption
entirely (forfeiting the monorepo's collaborative benefits). Neither
outcome is acceptable in a public, forkable template repository intended
to demonstrate best practices \citep{Wilson2014best}.

\subsection{Contribution}\label{contribution}

This paper introduces \breaktt{template\_pools\_rules\_tools}, a
\textbf{meta-project exemplar} that resolves this dilemma with a
four-module architecture (fig.~\ref{fig:architecture}). Each module
handles one resource category plus a fourth orchestration module:

{\def\LTcaptype{none} % do not increment counter
\begin{longtable}[]{@{}
  >{\raggedright\arraybackslash}p{(\linewidth - 4\tabcolsep) * \real{0.3333}}
  >{\raggedright\arraybackslash}p{(\linewidth - 4\tabcolsep) * \real{0.3333}}
  >{\raggedright\arraybackslash}p{(\linewidth - 4\tabcolsep) * \real{0.3333}}@{}}
\toprule\noalign{}
\begin{minipage}[b]{\linewidth}\raggedright
Module
\end{minipage} & \begin{minipage}[b]{\linewidth}\raggedright
Resource category
\end{minipage} & \begin{minipage}[b]{\linewidth}\raggedright
Key function
\end{minipage} \\
\midrule\noalign{}
\endhead
\bottomrule\noalign{}
\endlastfoot
\breaktt{src/fonds\_reader.py} & Data pools &
\breaktt{read\_all\_fonds()} \\
\breaktt{src/rules\_applier.py} & Governance rules &
\breaktt{validate\_against\_rules()} \\
\breaktt{src/tools\_invoker.py} & Executable tools &
\breaktt{discover\_tools()} \\
\breaktt{src/integration.py} & All three &
\breaktt{run\_integration\_demo()} \\
\end{longtable}
}

The architecture obeys three design invariants:

\begin{itemize}
\tightlist
\item
  \textbf{Read-only resource access}: no module writes to
  \texttt{fonds/}, \texttt{rules/}, or \texttt{tools/}. The Layer
  Contract in \texttt{AGENTS.md} enforces this at code-review time.
\item
  \textbf{Repo-root-relative discovery}: all path resolution uses
  \texttt{pathlib.Path(\_\_file\_\_).resolve().parents{[}N{]}} so that
  scripts work from any working directory.
\item
  \textbf{Graceful degradation}: every reader catches
  \breaktt{FileNotFoundError} and \texttt{yaml.YAMLError}, logs a
  warning, and returns a safe empty value. The pipeline never raises on
  a missing resource.
\end{itemize}

\subsection{Paper Organisation}\label{paper-organisation}

The remainder of this paper is structured as follows.
sec.~\ref{sec:pools} describes the fond layer and the
\texttt{fonds\_reader} module. sec.~\ref{sec:rules} describes the rules
layer and the \texttt{rules\_applier} module. sec.~\ref{sec:tools}
describes the tool layer and the \texttt{tools\_invoker} module.
sec.~\ref{sec:integration} presents the unified integration pipeline,
the manuscript variable token system, and a discussion of resilience
design. sec.~\ref{sec:conclusion} summarises key findings and future
directions.

The architecture overview in fig.~\ref{fig:architecture} provides a
visual map of these relationships. Runtime statistics collected during
integration are visualised in fig.~\ref{fig:counts}.

\newpage

\section{Pools: Fonds as Passive Data Resources}\label{sec:pools}

\subsection{What Is a Fond?}\label{what-is-a-fond}

A \textbf{fond} is a versioned, read-only data pool that any project in
the repository can consume without modifying. The term evokes the
culinary concept of a concentrated stock --- a stable base that enriches
whatever is built on top of it. Fonds live under the top-level
\texttt{fonds/\textless{}scope\textgreater{}/\textless{}name\textgreater{}/}
directory, each containing a manifest file (\texttt{fonds.yaml}), a
\texttt{data/} subdirectory, and optional documentation. This
architecture separates \emph{data ownership} from \emph{data usage}:
research projects in \texttt{projects/} are consumers, not producers, of
fond data. The separation prevents the accretion of project-specific
mutations in shared resources --- a recurring source of reproducibility
failures in collaborative research software \citep{Wilson2014best}.

The three-layer taxonomy (bibliography, contacts, datasets) maps to the
three most common categories of curated research data: citable
literature, human collaboration networks, and input/output datasets.
Each category carries its own schema enforced by
\breaktt{rules/templates/template\_manuscript\_rules}.

\subsection{The Three Template Fonds}\label{the-three-template-fonds}

\subsubsection{\texorpdfstring{\breaktt{template\_bibliography}}{template_bibliography}}\label{template_bibliography}

The \breaktt{template\_bibliography} fond is a curated reference library
stored in two formats: a BibTeX file (\breaktt{data/references.bib}) as
source of truth, and a flat CSV export (\breaktt{data/references.csv})
for programmatic querying. Deduplication is enforced on the cite key
(the primary CSV column). The collection spans foundational
machine-learning works --- the transformer architecture
\citep{Vaswani2017attention}, early convolutional network research
\citep{LeCun1998gradient}, and large-scale language model pre-training
\citep{Brown2020gpt3} --- alongside software-engineering references on
best practices \citep{Wilson2014best} and robust software design
\citep{Taschuk2017ten}. In the current integration run, the fond
contains \textbf{\{\{integration.bib\_entries\}\} entries}.

The bibliography fond illustrates the \emph{registry pattern}
\citep{Fowler2002patterns}: a single authoritative list is maintained
centrally, and all projects reference it rather than keeping private
copies. This guarantees citation consistency across all exemplar
manuscripts.

\subsubsection{\texorpdfstring{\breaktt{template\_contacts}}{template_contacts}}\label{template_contacts}

The \breaktt{template\_contacts} fond holds a registry of research
collaborators, advisors, and reviewers. Each entry is a YAML object with
required fields \texttt{id} (a unique slug), \texttt{name}, and
\texttt{email}, plus optional fields \texttt{affiliation},
\texttt{role}, \texttt{orcid}, \texttt{website}, and \texttt{notes}. The
YAML file (\breaktt{data/contacts.yaml}) is the source of truth; a JSON
mirror (\breaktt{data/contacts.json}) supports consumers that prefer JSON
deserialization. Deduplication is enforced on the \texttt{id} field at
validation time.

\subsubsection{\texorpdfstring{\breaktt{template\_datasets}}{template_datasets}}\label{template_datasets}

The \breaktt{template\_datasets} fond catalogs dataset metadata:
provenance, licensing, format, size, access URLs, and research tasks. It
intentionally stores \emph{metadata only} --- no actual data binaries
are committed to the repository. This design aligns with the principle
that version control systems should track configuration and metadata
rather than large binary artefacts \citep{Kluyver2016jupyter}. Dataset
entries require \texttt{id}, \texttt{name}, \texttt{version}, and
\texttt{license} fields. Exemplar entries reference classic benchmarks
such as MNIST (introduced in \citep{LeCun1998gradient}) and large-scale
corpora used in language-model research \citep{Brown2020gpt3}.

\subsection{\texorpdfstring{The \texttt{fonds.yaml}
Manifest}{The fonds.yaml Manifest}}\label{the-fonds.yaml-manifest}

Every fond root must contain a \texttt{fonds.yaml} manifest with at
minimum three fields:

\begin{Shaded}
\begin{Highlighting}[]
\FunctionTok{type}\KeywordTok{:}\AttributeTok{ bibliography}\CommentTok{   \# bibliography | contacts | datasets}
\FunctionTok{description}\KeywordTok{:}\AttributeTok{ }\StringTok{"Human{-}readable description of the fond"}
\FunctionTok{version}\KeywordTok{:}\AttributeTok{ }\StringTok{"1.0"}
\FunctionTok{tags}\KeywordTok{:}\AttributeTok{ }\KeywordTok{[}\AttributeTok{curated}\KeywordTok{,}\AttributeTok{ exemplar}\KeywordTok{]}
\end{Highlighting}
\end{Shaded}

The \texttt{type} field governs which reader function is appropriate and
what schema the \texttt{data/} directory is expected to follow. The
\texttt{version} field is incremented whenever the schema changes,
enabling consumers to detect and handle schema drift without silent
failures.

\subsection{\texorpdfstring{The \texttt{fonds\_reader}
Module}{The fonds_reader Module}}\label{the-fonds_reader-module}

The \breaktt{src/fonds\_reader.py} module provides three reader functions
--- one per fond type --- plus a convenience aggregator:

\begin{Shaded}
\begin{Highlighting}[]
\ImportTok{from}\NormalTok{ src.fonds\_reader }\ImportTok{import}\NormalTok{ (}
\NormalTok{    read\_bibliography\_fond,}
\NormalTok{    read\_contacts\_fond,}
\NormalTok{    read\_datasets\_fond,}
\NormalTok{    read\_all\_fonds,}
\NormalTok{)}

\NormalTok{bib      }\OperatorTok{=}\NormalTok{ read\_bibliography\_fond()   }\CommentTok{\# dict | None}
\NormalTok{contacts }\OperatorTok{=}\NormalTok{ read\_contacts\_fond()       }\CommentTok{\# dict | None}
\NormalTok{datasets }\OperatorTok{=}\NormalTok{ read\_datasets\_fond()       }\CommentTok{\# dict | None}
\NormalTok{all\_fonds }\OperatorTok{=}\NormalTok{ read\_all\_fonds()          }\CommentTok{\# \{"bibliography": ..., "contacts": ..., "datasets": ...\}}
\end{Highlighting}
\end{Shaded}

Each reader resolves the repository root from
\texttt{pathlib.Path(\_\_file\_\_).resolve().parents{[}4{]}} and uses a
\texttt{try/except} block that catches both \breaktt{FileNotFoundError}
and \texttt{yaml.YAMLError}. Returning \texttt{None} on failure rather
than raising ensures that the integration pipeline degrades gracefully
when a fond has not yet been populated by a parallel agent
\citep{Taschuk2017ten}. In the current run,
\{\{integration.fonds\_loaded\}\} of 3 expected fonds were successfully
loaded (see fig.~\ref{fig:counts}).

\subsection{Resilience by Design}\label{resilience-by-design}

The fond layer enforces resilience at two levels. At the
\textbf{structural} level, readers tolerate missing fonds entirely. At
the \textbf{schema} level, the manifest version field allows consumers
to check compatibility before processing data. This two-level approach
means a fond can evolve its schema without breaking existing consumers
that have not yet been updated: the consumer detects the version
mismatch and either adapts or skips, rather than crashing silently on
malformed data.

\newpage

\section{Rules: Soft and Strong Governance}\label{sec:rules}

\subsection{The Role of Governance
Rules}\label{the-role-of-governance-rules}

Research software projects make dozens of implicit governance decisions:
what test-coverage threshold is acceptable, how manuscript sections
should be ordered, which citation fields are mandatory. Left implicit,
these decisions drift silently across projects in a monorepo, eroding
the consistency that makes the repository valuable as a public exemplar.
The rules layer makes governance explicit, versioned, and
machine-enforceable.

A \textbf{rule set} in the template repository is a directory under
\texttt{rules/\textless{}scope\textgreater{}/\textless{}name\textgreater{}/}
containing a typed manifest (\texttt{rules.yaml}) and two subdirectories
of rule files:

\begin{verbatim}
<name>/
├── rules.yaml       — manifest (type, scope, version, rule_kinds)
├── soft/            — Markdown guideline files (human-readable, prompt-like)
└── strong/          — YAML constraint schemas (machine-enforceable)
\end{verbatim}

This two-tier architecture reflects the distinction between
\emph{guidance} (which humans follow approximately) and
\emph{constraints} (which pipelines enforce precisely) --- a distinction
also recognised in enterprise application architecture
\citep{Fowler2002patterns}.

\subsection{Soft Rules: Style and Process
Guidelines}\label{soft-rules-style-and-process-guidelines}

Soft rules are Markdown files in \texttt{soft/}. They encode preferences
and conventions that cannot easily be expressed as boolean constraints
but that human reviewers and AI agents can apply contextually. Examples
include:

\begin{itemize}
\tightlist
\item
  \textbf{Style preferences}: ``Prefer active voice in manuscript
  sections.'' ``Use \texttt{\textbackslash{}module\{\}} macros for all
  code identifiers.''
\item
  \textbf{Process guidelines}: ``Tag pull requests with a review-stage
  label before requesting review.'' ``Update \texttt{TODO.md} before
  closing an issue.''
\item
  \textbf{Citation conventions}: ``Cite primary sources rather than
  textbooks where possible.''
\end{itemize}

Soft rules are treated as guidance: deviations surface as suggestions in
code review and manuscript audit reports, not as pipeline blockers. This
makes the soft layer suitable for evolving preferences that should not
break automated checks.

\subsection{Strong Rules: Hard
Constraints}\label{strong-rules-hard-constraints}

Strong rules are YAML files in \texttt{strong/}. Each file defines one
named constraint:

\begin{Shaded}
\begin{Highlighting}[]
\FunctionTok{rule}\KeywordTok{:}
\AttributeTok{  }\FunctionTok{name}\KeywordTok{:}\AttributeTok{ coverage{-}gate}
\AttributeTok{  }\FunctionTok{kind}\KeywordTok{:}\AttributeTok{ strong}
\AttributeTok{  }\FunctionTok{description}\KeywordTok{:}\AttributeTok{ }\StringTok{"Minimum test coverage threshold for src/ modules."}
\AttributeTok{  }\FunctionTok{applies\_to}\KeywordTok{:}\AttributeTok{ }\StringTok{"projects/*/src/"}
\AttributeTok{  }\FunctionTok{enforcement}\KeywordTok{:}\AttributeTok{ fail\_on\_violation}
\AttributeTok{  }\FunctionTok{constraints}\KeywordTok{:}
\AttributeTok{    }\FunctionTok{minimum\_line\_coverage}\KeywordTok{:}\AttributeTok{ }\DecValTok{90}
\AttributeTok{    }\FunctionTok{minimum\_branch\_coverage}\KeywordTok{:}\AttributeTok{ }\DecValTok{80}
\end{Highlighting}
\end{Shaded}

The \breaktt{enforcement:\ fail\_on\_violation} field signals that a
pipeline must halt and report when this rule is violated. Strong rules
are suitable for invariants that, if broken, indicate a genuine defect
rather than a style preference: coverage below 90\% means tests are
missing; a manuscript section without an abstract means the document is
incomplete.

\subsection{The Two Template Rule
Sets}\label{the-two-template-rule-sets}

\subsubsection{\texorpdfstring{\breaktt{template\_project\_rules}}{template_project_rules}}\label{template_project_rules}

This rule set governs software projects throughout the template
repository. Its strong rules currently comprise:

{\def\LTcaptype{none} % do not increment counter
\begin{longtable}[]{@{}
  >{\raggedright\arraybackslash}p{(\linewidth - 2\tabcolsep) * \real{0.5000}}
  >{\raggedright\arraybackslash}p{(\linewidth - 2\tabcolsep) * \real{0.5000}}@{}}
\toprule\noalign{}
\begin{minipage}[b]{\linewidth}\raggedright
File
\end{minipage} & \begin{minipage}[b]{\linewidth}\raggedright
Constraint
\end{minipage} \\
\midrule\noalign{}
\endhead
\bottomrule\noalign{}
\endlastfoot
\breaktt{strong/coverage-gate.yaml} & Minimum line coverage 90\%, branch
coverage 80\% for \texttt{src/} \\
\breaktt{strong/module-structure.yaml} & Required directory layout:
\texttt{src/}, \texttt{tests/}, \texttt{scripts/},
\texttt{manuscript/} \\
\end{longtable}
}

Its soft rules provide guidance on code style, commit message
conventions, and pull-request labelling.

\subsubsection{\texorpdfstring{\breaktt{template\_manuscript\_rules}}{template_manuscript_rules}}\label{template_manuscript_rules}

This rule set governs research manuscripts. Its strong rules comprise:

{\def\LTcaptype{none} % do not increment counter
\begin{longtable}[]{@{}
  >{\raggedright\arraybackslash}p{(\linewidth - 2\tabcolsep) * \real{0.5000}}
  >{\raggedright\arraybackslash}p{(\linewidth - 2\tabcolsep) * \real{0.5000}}@{}}
\toprule\noalign{}
\begin{minipage}[b]{\linewidth}\raggedright
File
\end{minipage} & \begin{minipage}[b]{\linewidth}\raggedright
Constraint
\end{minipage} \\
\midrule\noalign{}
\endhead
\bottomrule\noalign{}
\endlastfoot
\breaktt{strong/reference-schema.yaml} & Required BibTeX fields and
cite-key format constraints \\
\breaktt{strong/section-schema.yaml} & Required manuscript sections,
ordering, and minimum word counts \\
\end{longtable}
}

In the current pipeline run, \textbf{\{\{integration.rules\_sets\_ok\}\}
of 2 rule sets} validated successfully (fig.~\ref{fig:counts}).

\subsection{\texorpdfstring{The \texttt{rules\_applier}
Module}{The rules_applier Module}}\label{the-rules_applier-module}

The \breaktt{src/rules\_applier.py} module exposes three functions:

\begin{Shaded}
\begin{Highlighting}[]
\ImportTok{from}\NormalTok{ src.rules\_applier }\ImportTok{import}\NormalTok{ (}
\NormalTok{    load\_soft\_rules,}
\NormalTok{    load\_strong\_rules,}
\NormalTok{    validate\_against\_rules,}
\NormalTok{)}

\NormalTok{soft   }\OperatorTok{=}\NormalTok{ load\_soft\_rules(}\StringTok{"template\_project\_rules"}\NormalTok{)    }\CommentTok{\# list[dict]}
\NormalTok{strong }\OperatorTok{=}\NormalTok{ load\_strong\_rules(}\StringTok{"template\_project\_rules"}\NormalTok{)  }\CommentTok{\# list[dict]}
\NormalTok{result }\OperatorTok{=}\NormalTok{ validate\_against\_rules(}\StringTok{"template\_project\_rules"}\NormalTok{)}
\CommentTok{\# → \{"status": "ok" | "partial" | "missing", "soft\_count": N, "strong\_count": N\}}
\end{Highlighting}
\end{Shaded}

\breaktt{validate\_against\_rules()} performs two checks: (1) the
\texttt{rules.yaml} manifest is parseable YAML; (2) every rule file in
\texttt{soft/} and \texttt{strong/} is parseable YAML. It returns a
status of \texttt{"ok"} when both checks pass, \texttt{"partial"} when
the manifest exists but some rule files are missing or malformed, and
\texttt{"missing"} when the rule set directory is absent entirely. This
graduated status enables the integration pipeline to distinguish between
a rule set that has not yet been created (acceptable during active
development) and one that is present but broken (actionable defect).

\subsection{Rules and Manuscript
Variables}\label{rules-and-manuscript-variables}

Strong rule validation counts are injected into the manuscript through
the token system. The token \texttt{\{\{integration.rules\_sets\_ok\}\}}
expands to the count of rule sets that returned \texttt{status="ok"}
during the integration run. This creates a verifiable link between the
pipeline's actual behaviour and the manuscript's claims --- the
manuscript cannot assert successful validation without the pipeline
having actually succeeded.

\newpage

\section{Tools: Executable Entry Points}\label{sec:tools}

\subsection{What Is a Tool?}\label{what-is-a-tool}

A \textbf{tool} in the template repository is a directory under
\texttt{tools/\textless{}scope\textgreater{}/\textless{}name\textgreater{}/}
that packages one or more executable scripts behind a typed manifest
(\texttt{tools.yaml}). Tools provide the third layer of the resource
architecture: where fonds supply static data and rules supply governance
constraints, tools supply \emph{behaviour} --- computations, validation
runs, and agent skill invocations that projects can trigger without
re-implementing the underlying logic.

The tools layer deliberately mirrors the Unix philosophy of small,
composable utilities that communicate through standard interfaces
\citep{Raymond2003art}. Each tool declares its entrypoints (shell
scripts), its invocation contract (stdin/stdout/exit-code semantics),
and its capabilities (type, version, tags) in a single manifest file.
Consumers invoke tools through the \texttt{tools\_invoker} module
without needing to understand the tool's implementation details --- a
textbook application of the Facade pattern \citep{Gamma1994design}.

\subsection{\texorpdfstring{The \texttt{tools.yaml}
Manifest}{The tools.yaml Manifest}}\label{the-tools.yaml-manifest}

Every tool root must contain a \texttt{tools.yaml} manifest with the
following fields:

\begin{Shaded}
\begin{Highlighting}[]
\FunctionTok{type}\KeywordTok{:}\AttributeTok{ code\_executor | validator | skill | agent | renderer}
\FunctionTok{description}\KeywordTok{:}\AttributeTok{ }\StringTok{"Human{-}readable description of the tool"}
\FunctionTok{version}\KeywordTok{:}\AttributeTok{ }\StringTok{"1.0.0"}
\FunctionTok{tags}\KeywordTok{:}\AttributeTok{ }\KeywordTok{[}\AttributeTok{curated}\KeywordTok{,}\AttributeTok{ exemplar}\KeywordTok{,}\AttributeTok{ production}\KeywordTok{,}\AttributeTok{ experimental}\KeywordTok{]}
\FunctionTok{creator}\KeywordTok{:}\AttributeTok{ }\StringTok{"org/repo"}
\FunctionTok{license}\KeywordTok{:}\AttributeTok{ }\StringTok{"Apache{-}2.0"}
\FunctionTok{entrypoints}\KeywordTok{:}
\AttributeTok{  }\KeywordTok{{-}}\AttributeTok{ scripts/run.sh}
\AttributeTok{  }\KeywordTok{{-}}\AttributeTok{ scripts/validate.sh}
\end{Highlighting}
\end{Shaded}

The \texttt{type} field determines the invocation contract the consumer
should expect. The \texttt{entrypoints} list names the scripts that must
exist on disk; the \texttt{tools\_invoker} module validates their
presence at discovery time rather than at invocation time, making
failures visible early in the pipeline rather than at runtime.

\subsection{The Three Template Tools}\label{the-three-template-tools}

\subsubsection{\texorpdfstring{\breaktt{template\_code\_executor}}{template_code_executor}}\label{template_code_executor}

A generic code execution tool that accepts a JSON payload on standard
input and returns execution results as JSON. The invocation contract is:

{\def\LTcaptype{none} % do not increment counter
\begin{longtable}[]{@{}
  >{\raggedright\arraybackslash}p{(\linewidth - 6\tabcolsep) * \real{0.2500}}
  >{\raggedright\arraybackslash}p{(\linewidth - 6\tabcolsep) * \real{0.2500}}
  >{\raggedright\arraybackslash}p{(\linewidth - 6\tabcolsep) * \real{0.2500}}
  >{\raggedright\arraybackslash}p{(\linewidth - 6\tabcolsep) * \real{0.2500}}@{}}
\toprule\noalign{}
\begin{minipage}[b]{\linewidth}\raggedright
Entrypoint
\end{minipage} & \begin{minipage}[b]{\linewidth}\raggedright
stdin
\end{minipage} & \begin{minipage}[b]{\linewidth}\raggedright
stdout
\end{minipage} & \begin{minipage}[b]{\linewidth}\raggedright
exit code
\end{minipage} \\
\midrule\noalign{}
\endhead
\bottomrule\noalign{}
\endlastfoot
\texttt{scripts/run.sh} & \texttt{\{"code":\ str,\ "language":\ str\}} &
\texttt{\{"exit\_code":\ int,\ "stdout":\ str,\ "stderr":\ str\}} & 0 =
success \\
\breaktt{scripts/validate.sh} & --- & Human-readable validation report &
0 = valid \\
\end{longtable}
}

The code executor exemplifies tools that wrap a computational
capability. The JSON-in/JSON-out contract makes it easily composable
with pipeline orchestrators and agent frameworks.

\subsubsection{\texorpdfstring{\breaktt{template\_validator}}{template_validator}}\label{template_validator}

A JSON Schema validation tool. It reads a target document and a schema
from disk and reports validation results in human-readable form. The
entrypoint \breaktt{scripts/validate.sh} exits 0 when the document is
valid and non-zero with a detailed error message otherwise. The
validator tool is used in the project pipeline to validate
\breaktt{manuscript\_variables.json} against its expected schema before
manuscript rendering.

\subsubsection{\texorpdfstring{\texttt{template\_skill}}{template_skill}}\label{template_skill}

An agent skill invocation tool that wraps a Hermes-compatible skill
definition. The entrypoint \breaktt{scripts/invoke.sh} accepts a prompt
string on standard input and returns the agent response as text. This
tool type bridges the repository's tool architecture with external agent
frameworks, demonstrating that the same manifest-and-entrypoint pattern
applies equally to computational tools and AI agents.

\subsection{\texorpdfstring{The \texttt{tools\_invoker}
Module}{The tools_invoker Module}}\label{the-tools_invoker-module}

The \breaktt{src/tools\_invoker.py} module provides three public
functions:

\begin{Shaded}
\begin{Highlighting}[]
\ImportTok{from}\NormalTok{ src.tools\_invoker }\ImportTok{import}\NormalTok{ (}
\NormalTok{    discover\_tools,}
\NormalTok{    get\_tool\_entrypoints,}
\NormalTok{    validate\_tool\_scripts\_exist,}
\NormalTok{)}

\NormalTok{tools }\OperatorTok{=}\NormalTok{ discover\_tools()}
\CommentTok{\# → [\{"name": "template\_code\_executor", "manifest": \{...\}\}, ...]}

\NormalTok{eps }\OperatorTok{=}\NormalTok{ get\_tool\_entrypoints(}\StringTok{"template\_code\_executor"}\NormalTok{)}
\CommentTok{\# → ["scripts/run.sh", "scripts/validate.sh"]}

\NormalTok{result }\OperatorTok{=}\NormalTok{ validate\_tool\_scripts\_exist(}\StringTok{"template\_code\_executor"}\NormalTok{)}
\CommentTok{\# → \{"status": "ok" | "partial" | "missing", "missing\_scripts": [...]\}}
\end{Highlighting}
\end{Shaded}

\breaktt{discover\_tools()} scans \breaktt{tools/templates/} for
subdirectories containing a parseable \texttt{tools.yaml} and returns a
list of discovery records. It silently skips directories without a
manifest rather than raising, ensuring that partially-completed tool
directories do not block the pipeline.

\breaktt{validate\_tool\_scripts\_exist()} iterates over the manifest's
\texttt{entrypoints} list and checks each path against the filesystem.
It returns a structured result distinguishing between tools that are
fully ready (\texttt{"ok"}), partially configured (\texttt{"partial"}
--- some scripts missing), and entirely absent (\texttt{"missing"}). In
the current integration run,
\textbf{\{\{integration.tools\_discovered\}\} tools} were discovered
(fig.~\ref{fig:counts}), all with valid manifests.

\subsection{Tool Discovery and
Reproducibility}\label{tool-discovery-and-reproducibility}

The tools layer contributes to reproducibility by making the
\emph{availability} of computational capabilities explicit and
checkable. A project that hard-codes a path to a tool script becomes
brittle when the repository is reorganised. By contrast, a project that
calls \breaktt{discover\_tools()} and checks
\breaktt{validate\_tool\_scripts\_exist()} will detect missing
entrypoints at pipeline initialisation time and report them clearly,
rather than failing silently at execution time
\citep{Stodden2016enhancing}. This shift from implicit to explicit
dependency declaration is a key design principle of the template
architecture (see fig.~\ref{fig:architecture}).

\newpage

\section{Integration: Unified Pipeline and Token
Injection}\label{sec:integration}

\subsection{Architecture Overview}\label{architecture-overview}

The three resource layers described in sec.~\ref{sec:pools},
sec.~\ref{sec:rules}, and sec.~\ref{sec:tools} are orchestrated by a
single function in \breaktt{src/integration.py}. The
\breaktt{run\_integration\_demo()} function calls all three subsystems in
a defined order, collects their results into a structured dictionary,
and writes summary counts to
\breaktt{output/data/manuscript\_variables.json} for injection into this
manuscript at render time. fig.~\ref{fig:architecture} illustrates the
complete architecture.

\begin{verbatim}
run_integration_demo()
  ├── read_bibliography_fond()       → {"manifest": ..., "bibtex": ..., "csv_rows": [...]}
  ├── read_contacts_fond()           → {"manifest": ..., "contacts": [...]}
  ├── read_datasets_fond()           → {"manifest": ..., "datasets": [...]}
  ├── validate_against_rules("template_project_rules")    → {"status": "ok", ...}
  ├── validate_against_rules("template_manuscript_rules") → {"status": "ok", ...}
  ├── discover_tools()               → [{"name": ..., "manifest": ...}, ...]
  └── validate_tool_scripts_exist(<each tool>)            → {"status": "ok", ...}
\end{verbatim}

The function returns a top-level dict with keys \texttt{fonds},
\texttt{rules}, \texttt{tools}, and \texttt{summary}. The
\texttt{summary} sub-dict provides the counts that populate manuscript
tokens.

\subsection{Manuscript Variable
Tokens}\label{manuscript-variable-tokens}

The token injection system bridges the integration pipeline and the
manuscript prose. Tokens use double-brace syntax:
\texttt{\{\{integration.fonds\_loaded\}\}} expands to the integer count
of fonds successfully loaded during the most recent integration run.
Tokens are resolved by \breaktt{scripts/03\_generate\_manuscript.py},
which reads \breaktt{output/data/manuscript\_variables.json} and
substitutes each token before the manuscript is passed to the rendering
engine.

The current \breaktt{manuscript\_variables.json} contains the following
summary values (see fig.~\ref{fig:counts} for a visual representation):

{\def\LTcaptype{none} % do not increment counter
\begin{longtable}[]{@{}
  >{\raggedright\arraybackslash}p{(\linewidth - 2\tabcolsep) * \real{0.5000}}
  >{\raggedright\arraybackslash}p{(\linewidth - 2\tabcolsep) * \real{0.5000}}@{}}
\toprule\noalign{}
\begin{minipage}[b]{\linewidth}\raggedright
Token
\end{minipage} & \begin{minipage}[b]{\linewidth}\raggedright
Value
\end{minipage} \\
\midrule\noalign{}
\endhead
\bottomrule\noalign{}
\endlastfoot
\texttt{\{\{integration.fonds\_loaded\}\}} &
\{\{integration.fonds\_loaded\}\} \\
\texttt{\{\{integration.rules\_sets\_ok\}\}} &
\{\{integration.rules\_sets\_ok\}\} \\
\texttt{\{\{integration.tools\_discovered\}\}} &
\{\{integration.tools\_discovered\}\} \\
\texttt{\{\{integration.bib\_entries\}\}} &
\{\{integration.bib\_entries\}\} \\
\end{longtable}
}

This table is itself token-injected: the values shown are those produced
by the pipeline, not hard-coded by the manuscript author. If the
pipeline results change --- for example, because a new fond is added ---
re-running \breaktt{scripts/03\_generate\_manuscript.py} updates the
manuscript automatically, without manual editing. This property is
central to reproducibility: the manuscript's quantitative claims are
always consistent with the code that generated them
\citep{Stodden2016enhancing}.

\subsection{Methods: The Script
Pipeline}\label{methods-the-script-pipeline}

Three thin orchestration scripts govern the integration workflow
(fig.~\ref{fig:pipeline}):

{\def\LTcaptype{none} % do not increment counter
\begin{longtable}[]{@{}
  >{\raggedright\arraybackslash}p{(\linewidth - 4\tabcolsep) * \real{0.3333}}
  >{\raggedright\arraybackslash}p{(\linewidth - 4\tabcolsep) * \real{0.3333}}
  >{\raggedright\arraybackslash}p{(\linewidth - 4\tabcolsep) * \real{0.3333}}@{}}
\toprule\noalign{}
\begin{minipage}[b]{\linewidth}\raggedright
Script
\end{minipage} & \begin{minipage}[b]{\linewidth}\raggedright
Purpose
\end{minipage} & \begin{minipage}[b]{\linewidth}\raggedright
Key output
\end{minipage} \\
\midrule\noalign{}
\endhead
\bottomrule\noalign{}
\endlastfoot
\breaktt{scripts/01\_validate\_sources.py} & Validate presence and
well-formedness of all resources & Console report \\
\breaktt{scripts/02\_run\_integration.py} & Run
\breaktt{run\_integration\_demo()} and print JSON summary & Console
JSON \\
\breaktt{scripts/03\_generate\_manuscript.py} & Write
\breaktt{output/data/manuscript\_variables.json} & JSON file \\
\end{longtable}
}

Each script is \texorpdfstring{\ensuremath{\leq}}{<=} 50 lines, imports all business logic from
\texttt{src/}, and uses \texttt{argparse} for optional flags. This
thin-orchestrator pattern \citep{Wilson2014best} ensures that all
testable logic is in \texttt{src/} at \texorpdfstring{\textgreater{}=}{>=} 90\% coverage, while the scripts
themselves remain readable without a test suite.

\subsection{Resilience Design}\label{resilience-design}

The integration layer enforces resilience at three levels, corresponding
to the three failure modes a monorepo integration pipeline encounters:

\begin{enumerate}
\def\labelenumi{\arabic{enumi}.}
\item
  \textbf{Resource absence}: A fond, rule set, or tool directory may not
  yet exist if the resource was created by a parallel agent that has not
  yet completed. All three reader modules return \texttt{None} or empty
  collections in this case, and \breaktt{run\_integration\_demo()}
  reports the missing resource in the \texttt{summary} dict without
  raising.
\item
  \textbf{Schema malformation}: A manifest may be present but contain
  invalid YAML or missing required fields. Readers catch
  \texttt{yaml.YAMLError} and return a degraded result
  (\breaktt{status="partial"} in rule validation; \texttt{None} in fond
  reading) rather than propagating the parse error.
\item
  \textbf{Script absence}: A tool may declare entrypoints in its
  manifest that have not yet been created.
  \breaktt{validate\_tool\_scripts\_exist()} detects this at discovery
  time and returns \breaktt{status="partial"} with a list of missing
  script paths, so the pipeline can report the defect without attempting
  to invoke a non-existent script.
\end{enumerate}

This three-level resilience design reflects the principle that
shared-resource pipelines should fail \emph{informatively} rather than
\emph{catastrophically} --- surfacing the cause of incompleteness in
structured output that downstream consumers can act on
\citep{Taschuk2017ten}.

\subsection{Test Coverage}\label{test-coverage}

The four \texttt{src/} modules are covered by 38 tests across four test
files in \texttt{tests/}. Tests use real file paths, real YAML files,
and real BibTeX content rather than mocks, ensuring that coverage
numbers reflect genuine code paths through the resource-discovery logic.
The current coverage report shows \texorpdfstring{\textgreater{}=}{>=} 90\% line coverage for all four
modules. The \breaktt{tests/test\_integration.py} suite includes an
end-to-end test that calls \breaktt{run\_integration\_demo()} and asserts
that the \texttt{summary} dict contains the expected keys with
non-negative integer values --- a contract test that verifies the token
injection pipeline's data source.

\newpage

\section{Conclusion}\label{sec:conclusion}

\subsection{Summary}\label{summary}

This paper has presented \breaktt{template\_pools\_rules\_tools}, a
meta-project exemplar demonstrating how research software projects
embedded in a monorepo can integrate three categories of shared
resources --- data pools (fonds), governance rules, and executable tools
--- without tight coupling to any specific resource instance. The key
contributions are:

\begin{enumerate}
\def\labelenumi{\arabic{enumi}.}
\item
  \textbf{A four-module architecture} (\texttt{fonds\_reader},
  \texttt{rules\_applier}, \texttt{tools\_invoker},
  \texttt{integration}) that provides a canonical pattern for
  resource-aware project code in the template repository. Each module is
  independently testable and independently deployable.
\item
  \textbf{A typed manifest convention} (\texttt{fonds.yaml},
  \texttt{rules.yaml}, \texttt{tools.yaml}) that makes resource
  capabilities explicit and checkable at pipeline initialisation time,
  shifting failure detection from runtime to startup --- a significant
  improvement for reproducibility \citep{Wilson2014best}.
\item
  \textbf{A token injection pipeline} that links manuscript prose to
  integration runtime statistics through \texttt{\{\{integration.*\}\}}
  tokens, ensuring that quantitative claims in the manuscript are always
  generated by the pipeline rather than authored manually. In the
  current run this covered \{\{integration.fonds\_loaded\}\} fonds,
  \{\{integration.rules\_sets\_ok\}\} rule sets,
  \{\{integration.tools\_discovered\}\} tools, and
  \{\{integration.bib\_entries\}\} bibliography entries.
\item
  \textbf{A three-level resilience design} --- resource absence, schema
  malformation, and script absence --- that allows the pipeline to
  degrade gracefully and report failures informatively rather than
  crashing, consistent with best practices for robust research software
  \citep{Taschuk2017ten}.
\end{enumerate}

\subsection{Design Decisions
Revisited}\label{design-decisions-revisited}

Several design decisions deserve emphasis as lessons for future exemplar
authors:

\textbf{Repo-root-relative discovery} is non-negotiable in a forkable
monorepo. Any hard-coded absolute path or working-directory assumption
will fail when the repository is cloned to a different location. The
\texttt{pathlib.Path(\_\_file\_\_).resolve().parents{[}N{]}} idiom used
throughout \texttt{src/} ensures that discovery works from any working
directory.

\textbf{Graceful fallbacks over exceptions} is the correct trade-off for
a resource-consumption layer. The alternative --- raising
\breaktt{FileNotFoundError} when a fond is missing --- would make the
integration pipeline fragile to the ordering of parallel agent writes.
Returning \texttt{None} and logging a warning allows the pipeline to
produce a complete, if partial, result dict that downstream consumers
can reason about.

\textbf{Real-file tests over mocks} is required for a template exemplar.
Mocks can pass tests while the real discovery logic is silently broken.
Tests that use real file paths exercise the full path-resolution chain
and catch regressions that mocks would miss.

\textbf{Thin orchestration scripts} keep the scripts directory readable
and focus all testable logic in \texttt{src/}. A script that does more
than parse arguments, call a \texttt{src/} function, and write output is
accumulating business logic that belongs elsewhere.

\subsection{Future Directions}\label{future-directions}

The three-layer architecture described here is deliberately extensible.
The most natural extension is a fourth resource category:
\textbf{models}
(\texttt{models/\textless{}scope\textgreater{}/\textless{}name\textgreater{}/})
for pre-trained machine learning models and their inference scripts. The
pattern would follow exactly the same manifest-and-reader design, with
\texttt{models.yaml} declaring type, version, and entrypoints, and a
\texttt{models\_loader} module providing \breaktt{discover\_models()} and
\breaktt{validate\_model\_files\_exist()}. Adding this layer to
\breaktt{template\_pools\_rules\_tools} would require only a new reader
module and a new section in the integration orchestrator --- the
existing architecture places no constraints on the number of resource
categories.

A second direction is \textbf{cross-fond validation}: checking that cite
keys used in a project's \breaktt{manuscript/references.bib} are a subset
of those available in \breaktt{fonds/templates/template\_bibliography/}.
This would require a small cross-reference function in the integration
module and would strengthen the reproducibility guarantee by ensuring
that manuscript citations always trace back to curated, shared sources.

The exemplar is ready to fork. A project team wishing to adopt this
architecture need only copy the four \texttt{src/} modules, update the
resource directory paths in each reader, and replace the exemplar
fond/rules/tool names with their own.



\clearpage

\bibliography{references}
\end{document}
